\section{Gestión de la Información (Cumplimiento ISO 26515)}
\label{sec:s3_gestion}

Las funcionalidades concurrentes exigen documentación operativa específica: manuales de monitoreo
(Sección~\ref{ssec:s3_gi_monitoreo}) y guías de resolución de conflictos
(Sección~\ref{ssec:s3_gi_troubleshooting}).\par

\subsection{Manuales de Monitoreo de Servicios en Tiempo Real}
\label{ssec:s3_gi_monitoreo}
Se generaron manuales operativos para el monitoreo de los servicios web en tiempo real (estado de
los canales Realtime, salud de la conexión y latencia), apoyados en Spring Boot Actuator y en el
\comando{HealthController}. La tabla~\ref{tab:s3_monitoreo} resume las señales monitorizadas.

\begin{table}[!ht]
    \centering
    \renewcommand{\arraystretch}{1.3}
    \fontsize{9}{10.8}\selectfont\sffamily
    \begin{tabularx}{\textwidth}{|l|X|}
        \hline
        \rowcolor{colorPrimario}
        \textcolor{white}{\textbf{Señal}} & \textcolor{white}{\textbf{Fuente / criterio}} \\ \hline
        Salud del servicio & \comando{/actuator/health} (UP/DOWN). \\ \hline
        \rowcolor{grisTecnico}
        Estado de canales & Suscripciones Realtime activas vs. esperadas. \\ \hline
        Latencia de IA & Tiempos de respuesta de Gemini y tasa de 429/5xx. \\ \hline
        \rowcolor{grisTecnico}
        Archivos huérfanos & Conteo purgado por el \textit{worker} asíncrono. \\ \hline
    \end{tabularx}
    \caption{Señales de monitoreo de los servicios en tiempo real (Sprint 3).}
    \label{tab:s3_monitoreo}
\end{table}

\subsection{Guías de Troubleshooting de Concurrencia}
\label{ssec:s3_gi_troubleshooting}
Se redactaron guías específicas para la resolución de conflictos derivados de la concurrencia,
recogidas en la tabla~\ref{tab:s3_troubleshooting}.

\begin{table}[!ht]
    \centering
    \renewcommand{\arraystretch}{1.3}
    \fontsize{9}{10.8}\selectfont\sffamily
    \begin{tabularx}{\textwidth}{|X|X|}
        \hline
        \rowcolor{colorPrimario}
        \textcolor{white}{\textbf{Síntoma}} & \textcolor{white}{\textbf{Mitigación}} \\ \hline
        Mensaje duplicado tras reconexión & Idempotencia por clave de deduplicación (V7). \\ \hline
        \rowcolor{grisTecnico}
        Condición de carrera en el pipeline & Mutaciones vía RPC \comando{SECURITY DEFINER} (V8). \\ \hline
        Canal Realtime caído & Reintentos de suscripción con \textit{backoff}. \\ \hline
        \rowcolor{grisTecnico}
        Saturación de IA (429) & Mapeo \comando{AiServiceException} y reintento diferido. \\ \hline
    \end{tabularx}
    \caption{Guía de troubleshooting de concurrencia (Sprint 3).}
    \label{tab:s3_troubleshooting}
\end{table}

La tabla Historias~vs.~Documentación cubrió las Épicas de proyectos, conexiones y reclutador.

\begin{TablaHSDoc}
    EPIC-04 & Vitrina de proyectos & Sí (S3) & Sí (S3) \\ \hline
    EPIC-05 & Conexiones externas & Sí (S3) & Sí (S3) \\ \hline
    EPIC-07 & Reclutador e IA & Sí (S3) & Sí (S3) \\ \hline
\end{TablaHSDoc}

\subsection{Productos de Información del Incremento (ISO 15289)}
\label{ssec:s3_gi_productos}
La tabla~\ref{tab:s3_productos} relaciona los productos de información de las funcionalidades
avanzadas con su artefacto de origen.

\begin{table}[!ht]
    \centering
    \renewcommand{\arraystretch}{1.3}
    \fontsize{9}{10.8}\selectfont\sffamily
    \begin{tabularx}{\textwidth}{|X|X|}
        \hline
        \rowcolor{colorPrimario}
        \textcolor{white}{\textbf{Producto de información}} & \textcolor{white}{\textbf{Artefacto de origen}} \\ \hline
        Manual de monitoreo Realtime & Actuator + \comando{HealthController} + canales. \\ \hline
        \rowcolor{grisTecnico}
        Guía de troubleshooting de concurrencia & RPCs V7/V8 e idempotencia. \\ \hline
        Documentación de RPCs del reclutador & Funciones \comando{SECURITY DEFINER} (V6/V11). \\ \hline
    \end{tabularx}
    \caption{Productos de información del Sprint 3 y su trazabilidad.}
    \label{tab:s3_productos}
\end{table}

\clearpage
